\emph{Setting.} Conventions are those of Definition \claimref{obj.synthesis.gtg-flat-spacetime-gauge-fields} and Fact \claimref{synthesis.imported-gtg-dictionary}: signature $(+,-,-,-)$, store curvature $\mathcal{R}(B) = -R(\hb(B))$, Ricci map $\mathcal{R}(b) = \partial_{a} \cdot \mathcal{R}(a \wedge b)$, scalar $\mathcal{R} = \partial_{a} \cdot \mathcal{R}(a)$, Einstein map $\mathcal{G}(a) = \mathcal{R}(a) - \tfrac12 a\,\mathcal{R}$, field equation $\mathcal{G}(a) + \Lambda a = -8\pi G\,T(a)$ (de Sitter has $\mathcal{R}(B) = \tfrac{\Lambda}{3}B$, $\mathcal{R} = 4\Lambda$), Schouten map $\Sch(a) = \tfrac12(\mathcal{R}(a) - \tfrac16 a\,\mathcal{R})$, and $\mathcal{R}(A \wedge B) = \mathcal{W}(A \wedge B) + \Sch(A) \wedge B + A \wedge \Sch(B)$. $D_{A}$ is the covariant derivative along $A$; for a map $S$ of vectors, $(D_{A}S)(B) = D_{A}(S(B)) - S(D_{A}B)$.

\emph{(i).} By Theorem \claimref{synthesis.gtg-tractor-gauge-field-reproduces-tractor-connection}(iv) the field strength $F_{T}$ of $\Omega_{T}$ has three blocks: the spacetime (Lorentz) block $-\mathcal{W}(A \wedge B)$, the $\nInf$ block the torsion, zero here, and the $\nO$ block the Cotton term $(D_{A}\Sch)(B) - (D_{B}\Sch)(A)$. So $F_{T} = 0$ implies $\mathcal{W} = 0$. Conversely let $\mathcal{W} = 0$ on a region. By Fact \claimref{correspondence.imported-weyls-theorem} about each point there is a gauge with $\hb^{-1}(a) = \phi\,a$. By Fact \claimref{correspondence.imported-tractor-connection-local-twistor-transport}(i) the tractor connection is a connection of the conformal class: a change of scale acts on $\Omega_{T}$ as a gauge transformation by a rotor of the parabolic subgroup, under which $F_{T}$ transforms by conjugation. In the flat scale $\phi = 1$ ($\Omega = 0$, $\Sch = 0$) one has $\Omega_{T}(a) = -\nInf \wedge a$, constant, so $F_{T}(a \wedge b) = \Omega_{T}(a) \times \Omega_{T}(b)$; since $\nInf$ is null and orthogonal to $a$ and $b$, $(\nInf \wedge a)(\nInf \wedge b) = \nInf a \nInf b = -\nInf^{2}ab = 0$ and likewise in the other order, so the commutator vanishes and $F_{T} = 0$; hence $F_{T} = 0$ in every scale of the class. The twistor attachment works with one conformal spinor space $\CS$ (Proposition \claimref{field-theory.twistor-attachment}). Part (ii) of the same Fact makes local twistor transport the spin action of $\Omega_{T}$ on $\CS$, a bivector $B$ acting by left multiplication $\psi \mapsto \tfrac12 B\psi$; this is the defining representation of $\su(2,2)$ on the four-dimensional space $\CS$, which is faithful (left multiplication by a nonzero bivector of $\Cl(2,4)$ is nonzero on $\CS$), so twistor transport is flat exactly where $F_{T} = 0$; a flat connection has trivial transport around every contractible loop, and a nonzero field strength gives a nontrivial transport around small loops in its plane at second order (the standard relation between curvature and holonomy, cited from memory). The points with $\mathcal{W} \neq 0$ are classified by Petrov type in the result \claimref{synthesis.attachment-transport-is-flat-only-where-weyl-vanishes}; only the flat case is used here. The check computes $F_{T}$ in all blocks at two rational points of two conformally flat fields with nonzero scalar curvature and finds it zero, and on a field that is not conformally flat finds $F_{T} \neq 0$ with spacetime block $-\mathcal{W}$. The formula $\mathcal{R} = 6\,\Box\phi/\phi^{3}$ is the standard conformal-change formula of the scalar curvature in the store's sign; the check computes $\mathcal{R}$ from the field strength by exact Taylor jets at four points of two fields and confirms it.

\emph{(ii).} From the field equation, taking the trace ($\partial_{a} \cdot a = 4$, $\partial_{a} \cdot \mathcal{G}(a) = -\mathcal{R}$), $\mathcal{R} = 4\Lambda + 8\pi G\tau$, so $\mathcal{R}(a) = -8\pi G\,(T(a) - \tfrac{\tau}{2}a) + \Lambda a$ and
\[
\Sch(a) = \tfrac12\bigl(-8\pi G\,(T(a) - \tfrac{\tau}{2}a) + \Lambda a - \tfrac16 a\,(4\Lambda + 8\pi G\tau)\bigr) = -4\pi G\,S(a) + \tfrac{\Lambda}{6}a, \qquad S(a) = T(a) - \tfrac{\tau}{3}a .
\]
By (i) the field is twistor-flat, so the Cotton block vanishes: $(D_{A}\Sch)(B) = (D_{B}\Sch)(A)$. Since $D_{A}\,\id = 0$ (the covariant derivative is metric and torsion-free) and $\Lambda$ is constant, this is $-4\pi G\,((D_{A}S)(B) - (D_{B}S)(A)) = 0$, the curl condition.

(a) In the index reading ($S_{bc} = w\,u_{b}u_{c} - \tfrac{\rho}{3}g_{bc}$ with $w = \rho + p$, which is $S$ for the perfect fluid since $\tau = \rho - 3p$), contract the curl $\nabla_{a}S_{bc} - \nabla_{b}S_{ac}$ with $u^{b}u^{c}$. Using $u^{c}\nabla_{a}u_{c} = 0$ (from $Q(u) = 1$), $u^{b}u^{c}\nabla_{a}S_{bc} = \nabla_{a}f$ with $f = \tfrac23\rho + p$, and $u^{b}u^{c}\nabla_{b}S_{ac} = u^{b}\nabla_{b}(f u_{a}) - S_{ac}\dot u^{c} = \dot f u_{a} + f\dot u_{a} + \tfrac{\rho}{3}\dot u_{a} = \dot f u_{a} + w\dot u_{a}$, where a dot is $u \cdot \nabla$. So
\[
u^{b}u^{c}(\nabla_{a}S_{bc} - \nabla_{b}S_{ac}) = \nabla_{a}f - u_{a}\dot f - w\,\dot u_{a} .
\]
This is an identity among first derivatives of covariant quantities, so it holds on every field once it holds at a point of the flat field for arbitrary first jets (at any point one may choose the gauge in which the connection vanishes there); the check verifies it exactly at a rational point for two sets of polynomial $\rho$, $p$ and a unit $u$ given by a rational parametrisation. Conservation of $T$, $\nabla^{b}T_{ab} = u_{a}\nabla^{b}(w u_{b}) + w\dot u_{a} - \nabla_{a}p = 0$, projected orthogonally to $u$, gives $w\dot u_{a} = \nabla_{a}p - u_{a}\dot p$. Substituting, the contraction equals $\tfrac23(\nabla_{a}\rho - u_{a}\dot\rho)$; the check verifies this combination as well. The curl condition makes it zero, so $\nabla\rho = u\,(u \cdot \nabla\rho)$. For static dust with varying density the curl is nonzero: with $u = e_{t}$, $p = 0$ on the flat field its contraction has spatial components $\tfrac23\nabla\rho$, as the check confirms. This agrees with Stephani's classification of conformally flat perfect fluids (cited from memory), whose densities are constant on the rest spaces.

(b) If $T = 0$ then $\Sch = \tfrac{\Lambda}{6}\id$ and $\mathcal{W} = 0$ give $\mathcal{R}(A \wedge B) = \tfrac{\Lambda}{6}(A \wedge B + A \wedge B) = \tfrac{\Lambda}{3}A \wedge B$; the check confirms this on the conformally flat de Sitter fields $\phi = 1/(1 - \Lambda Q(x)/12)$ for two values of $\Lambda$ at two points each. At $\Lambda = 0$ the field strength vanishes. The Schwarzschild field has vanishing Ricci map and $\mathcal{R}(\sigma_{r}) = \tfrac{2M}{r^{3}}\sigma_{r}$ (dictionary, store sign), so $\mathcal{W} = \mathcal{R} \neq 0$; a gravitational wave at $\Lambda = 0$ is a vacuum field of nonzero field strength, hence $\mathcal{W} \neq 0$ by (b). Neither is twistor-flat.

\emph{(iii).} On $\hb^{-1}(a) = \phi\,a$, $\mathcal{W} = 0$ identically by (i), and by (i) $\mathcal{R} = -24\pi G\tau$ reads $6\,\Box\phi/\phi^{3} = -24\pi G\tau$, that is $\Box\phi = -4\pi G\tau\phi^{3}$. Solving for $\partial_{t}^{2}\phi$, the Taylor coefficient of $\phi$ with time exponent $k + 2$ is fixed by coefficients of time exponent at most $k + 1$ (the spatial Laplacian keeps the time exponent, and the product $\tau\phi^{3}$ involves only coefficients of $\phi$ with time exponent at most $k$), so from Cauchy data every coefficient is determined, and no equation constrains $\tau$, since there is one equation for one unknown and no differential identity links it to the matter. The check carries this out to total order four for a trace depending on $t$, $x$, $y$, $z$, builds the field, computes its scalar curvature from the field strength as a jet, and finds $\mathcal{R} + 24\pi G\tau = 0$ through jet order two (with $4\pi G$ carried as a rational formal value), while a different trace fails; the field's tractor field strength is zero at a second point. For analytic data the formal solution converges (Cauchy-Kovalevskaya), and for smooth data local existence holds by the standard theory of semilinear wave equations (both cited from memory); only the formal statement is used for consistency. The Bianchi identities hold for the field strength of every gauge field. The matter is conserved by its own equations, as for any matter coupled through the derived metric (displacement-gauge invariance of the matter action); the gravitational equations add no condition, unlike Einstein's, whose contracted Bianchi identity requires $\nabla \cdot T = 0$ of the source. For a static weak field $\phi = 1 + \psi$, to first order $\Box\phi = -\nabla^{2}\psi$, so Nordstr\"om's equation gives $\nabla^{2}\psi = 4\pi G\rho$ (dust, $\tau = \rho$), and since $g_{tt} = \phi^{2} = 1 + 2\psi + \dots$ a slow test body moves in the Newtonian potential $\psi$. The trace of Einstein's equations, $\mathcal{R} = 4\Lambda + 8\pi G\tau$ from (ii), gives at $\Lambda = 0$ instead $-6\nabla^{2}\psi = 8\pi G\rho$, $\nabla^{2}\psi = -\tfrac{4\pi}{3}G\rho$. These equations are Nordstr\"om's theory in the form of Einstein and Fokker (Fact \claimref{synthesis.imported-nordstrom-scalar-theory-einstein-fokker-form}); the computation above re-derives what is used.

\emph{(iv).} By (iii), Nordstr\"om's equations on twistor-flat fields satisfy the Bianchi identities, are compatible with the conservation of every stress, are covariant (both equations are statements about the field strength and the stress trace), are of second order in $\phi$, and have the Newtonian static limit; by (ii) they are not Einstein's equations, which on this class have no solution with a perfect-fluid source of varying rest density, whereas Nordstr\"om's equations are solvable for every trace. So these consistency conditions do not imply Einstein's equations. Fact \claimref{synthesis.imported-lovelock-and-anderson-duchamp-natural-lagrangians}(a) gives the step that does: a natural symmetric map of the field, of second order and divergence-free for every field, is $c_{1}\mathcal{G}(a) + c_{2}a$; equating such a map to a multiple of $T(a)$ gives Einstein's equations with free $G$ and $\Lambda$. Nordstr\"om's equations are not of that form (the Weyl condition is unsourced and the scalar equation sees only $\tau$), so the clause that the whole stress map is the source is what separates the two, and it is not among the consistency conditions. Nordstr\"om's theory is excluded by observation (null geodesics are conformally invariant, so it deflects no light; its perihelion shift is minus one sixth of Einstein's; Fact \claimref{synthesis.imported-nordstrom-scalar-theory-einstein-fokker-form}): it serves here as a consistent counter-model, not as a candidate theory.
